% 本章使用的 BibTeX 键：hartleySturm1997。

\section{校正变换}
\label{sec:rectifying-transformation}

为了校正左图像，需要计算一个变换，将 PPM
$\widetilde{P}_{o1}=[Q_{o1}\mid\tilde{q}_{o1}]$ 的图像平面映射到
$\widetilde{P}_{n1}=[Q_{n1}\mid\tilde{q}_{n1}]$ 的图像平面。下面将看到，所求变换是由
$3\times3$ 矩阵
\[
  T_1=Q_{n1}Q_{o1}^{-1}
\]
给出的共线性变换。对右图像也有相同的结果。

对于任意三维点 $w$，有

\begin{equation}
  \begin{cases}
    \tilde{m}_{o1}\simeq\widetilde{P}_{o1}\tilde{w},\\
    \tilde{m}_{n1}\simeq\widetilde{P}_{n1}\tilde{w}
  \end{cases}
  \label{eq:original-new-projection}
\end{equation}

根据式~(8)，光线方程如下（因为立体校正不移动光心）：

\begin{equation}
  \begin{cases}
    w=c_1+\lambda_o Q_{o1}^{-1}\tilde{m}_{o1},\qquad \lambda_o\in\mathbb{R},\\
    w=c_1+\lambda_n Q_{n1}^{-1}\tilde{m}_{n1},\qquad \lambda_n\in\mathbb{R}
  \end{cases}
  \label{eq:original-new-rays}
\end{equation}

因此，

\begin{equation}
  \tilde{m}_{n1}=\lambda Q_{n1}Q_{o1}^{-1}\tilde{m}_{o1},\qquad \lambda\in\mathbb{R}
  \label{eq:rectifying-collinearity}
\end{equation}

随后将变换 $T_1$ 应用于原始左图像，以生成校正图像，如图~5 所示。注意，校正图像的
像素（整数坐标位置）一般对应于原始图像平面中的非整数位置。因此，校正图像的灰度值
通过双线性插值计算。

可以直接使用 $P_{n1}$、$P_{n2}$，通过三角测量从校正图像重建三维点
\citep{hartleySturm1997}。
